\documentclass[10pt]{article}
\usepackage[T1]{fontenc}
\usepackage[utf8]{inputenc}
\usepackage[french]{babel}
\usepackage{amsmath,amssymb,mathrsfs,mathpazo}
\usepackage{geometry}
\geometry{a4paper,margin=2cm}
\usepackage{graphicx,array,tabularx,tabularray,colortbl,multirow,multicol,booktabs,enumitem,pifont,ulem,fancybox,diagbox,cancel}
\usepackage{pstricks,pstricks-add,pst-plot,pst-tree,pst-node,pst-text,pst-func,pst-math,pst-bspline,pst-3dplot,pst-eucl,pst-solides3d}
\let\pscalculate\undefined
\usepackage{pst-bezier}
\usepackage{tikz,pgfplots,tkz-tab}
\usetikzlibrary{arrows,arrows.meta,patterns,calc,positioning,trees,decorations.markings,shapes.geometric,shapes.misc,shapes.arrows,intersections,decorations.pathmorphing,decorations.pathreplacing,backgrounds,shadows,babel}
\pgfplotsset{compat=1.18}
\usepgfplotslibrary{fillbetween}
\usepackage[np]{numprint}
\usepackage{listings,makecell,pas-tableur,twoopt}
\let\barre\undefined
\usepackage[pstricks]{tablvar}
\definecolor{darkgreen}{RGB}{0,100,0}
\providecommand{\red}{\color{red}}
\providecommand{\blue}{\color{blue}}
\providecommand{\green}{\color{green}}
\providecommand{\black}{\color{black}}
\providecommand{\white}{\color{white}}
\providecommand{\gray}{\color{gray}}
\providecommand{\cyan}{\color{cyan}}
\providecommand{\magenta}{\color{magenta}}
\providecommand{\yellow}{\color{yellow}}
\providecommand{\violet}{\color{violet}}
\providecommand{\orange}{\color{orange}}
\providecommand{\euro}{\texteuro}
\providecommand{\EUR}{\texteuro}
\providecommand{\ieme}{\textsuperscript{e}}
\providecommand{\degres}{\ensuremath{^\circ}}
\providecommand{\degree}{\ensuremath{^\circ}}
\providecommand{\vv}[1]{\overrightarrow{#1}}
\providecommand{\Coord}[1]{#1}
\providecommand{\decosix}{}
\providecommand{\starredbullet}{\ensuremath{\bullet}}
\pagestyle{empty}
\setlength{\parindent}{0pt}
\usepackage{adjustbox}\DeclareUnicodeCharacter{2236}{:}\DeclareUnicodeCharacter{2190}{\ensuremath{\leftarrow}}
\providecommand{\R}{\mathbb{R}}
\providecommand{\N}{\mathbb{N}}
\providecommand{\Z}{\mathbb{Z}}
\providecommand{\Q}{\mathbb{Q}}
\providecommand{\D}{\mathbb{D}}
\providecommand{\C}{\mathbb{C}}
\providecommand{\e}{{\mathrm{e}}}
\providecommand{\vect}[1]{\overrightarrow{#1}}
\providecommand{\vectt}[1]{\overrightarrow{\mathrm{#1}}}
\providecommand{\vvt}[1]{\overrightarrow{\mathrm{#1}}}
\providecommand{\np}[2][]{{#2~#1}}
\providecommand{\vv}[1]{{\overrightarrow{#1}}}
\renewcommand{\d}{{\mathrm{d}}}
\providecommand{\degres}{{^\circ}}
\providecommand{\degre}{{^\circ}}
\providecommand{\blue}{\color{blue}}
\providecommand{\red}{\color{red}}
\providecommand{\enskip}{\quad}
\providecommand{\llbracket}{\lbrack\!\lbrack}
\providecommand{\rrbracket}{\rbrack\!\rbrack}
\providecommand{\euro}{{\text{€}}}
\providecommand{\no}{{\text{n°}}}
\providecommand{\interval}[2]{{[#1;#2]}}
\providecommand{\intervFO}[2]{{[#1;#2[}}
\providecommand{\dotfill}{\ldots}

\providecommand{\ds}{\displaystyle}
\providecommand{\textding}[1]{#1}
\providecommand{\Oij}{\(\mathrm{O};\vec{\imath},\vec{\jmath}\)}
\providecommand{\Oijk}{\(\mathrm{O};\vec{\imath},\vec{\jmath},\vec{k}\)}
\providecommand{\Ouv}{\(\mathrm{O};\vec{u},\vec{v}\)}
\providecommand{\intervFF}[2]{\ensuremath{[#1;#2]}}
\renewcommand{\euro}{\ifmmode\text{\texteuro}\else\texteuro\fi}
\DeclareUnicodeCharacter{00B0}{\ifmmode^{\circ}\else\textdegree\fi}
\DeclareUnicodeCharacter{20AC}{\ifmmode\text{\texteuro}\else\texteuro\fi}
% Énoncés : Éducation nationale. Transcription et corrigés : APMEP, Denis Vergès et contributeurs. Découpage et classement : CoopMaths, Sébastien Lozano, Éric Elter et contributeurs. Adaptation de la présentation : Maths Explorer.
% CC BY-SA : https://coopmaths.fr/www/about/
% https://coopmaths.fr/alea/static/bac/2022/tex/bac_2022_05_sujet1_madagascar_2.tex
\begin{document}
\noindent\textit{APMEP / CoopMaths. CC BY-SA. Présentation adaptée par Maths Explorer.}\par\medskip
\medskip

\begin{tabularx}{\linewidth}{|X|}\hline
\textbf{Principaux domaines abordés:} Suites ; Fonctions, Fonction logarithme.\\ \hline
\end{tabularx}

\medskip

\emph{Cet exercice est un questionnaire à choix multiples.\\
Pour chacune des questions suivantes, une seule des quatre réponses proposées
est exacte.\\
Une réponse exacte rapporte un point. Une réponse fausse, une réponse multiple ou
l'absence de réponse à une question ne rapporte ni n'enlève de point.\\
Pour répondre, indiquer sur la copie le numéro de la question et la lettre de la réponse choisie. Aucune justification n'est demandée.}

\bigskip
\begin{enumerate}
\item Un récipient contenant initialement 1 litre d'eau est laissé au soleil.

Toutes les heures, le volume d'eau diminue de 15\,\%.

Au bout de quel nombre entier d'heures le volume d'eau devient-il inférieur à un quart de litre ?

\begin{center}
\begin{tabularx}{\linewidth}{*{4}{X}}
\textbf{a.~~} 2 heures &\textbf{b.~~}  8 heures .&
\textbf{c.~~} 9 heures&\textbf{d.~~} 13 heures 
\end{tabularx}
\end{center}

\item On considère la fonction $f$ définie sur l'intervalle $]0~;~+\infty[$ par $f(x) = 4\ln (3x)$.

Pour tout réel $x$ de l'intervalle $]0~;~+\infty[$ , on a :

\begin{center}
\begin{tabularx}{\linewidth}{*{2}{X}}
\textbf{a.~~} $f(2x) = f(x) + \ln (24)$&\textbf{b.~~} $f(2x) = f(x) + \ln (16)$\\
\textbf{c.~~} $f(2x) = \ln (2) + f(x)$& \textbf{d.~~} $f(2x) = 2f(x)$
\end{tabularx}
\end{center}
\item On considère la fonction $g$ définie sur l'intervalle $]1~;~+\infty[$ par : 

\[g(x) = \dfrac{\ln (x)}{x - 1}.\]


On note $\mathcal{C}_g$ la courbe représentative de la fonction $g$ dans un repère orthogonal. La courbe $\mathcal{C}_g$ admet :

\begin{center}
\begin{tabularx}{\linewidth}{*{2}{X}}
\textbf{a.~~} une asymptote verticale
et une asymptote horizontale.&\textbf{b.~~} une asymptote verticale
et aucune asymptote horizontale.\\
\textbf{c.~~} aucune asymptote verticale et une asymptote horizontale.&\textbf{d.~~} 
aucune asymptote verticale .et aucune asymptote horizontale.
\end{tabularx}
\end{center}
\end{enumerate}

Dans la suite de l'exercice, on considère la fonction $h$ définie sur l'intervalle ]0~;~2] par:

\[h(x) = x^2(1 + 2\ln (x)).\]

On note $\mathcal{C}_h$ la courbe représentative de $h$ dans un repère du plan. 

On admet que $h$ est deux fois dérivable sur l'intervalle ]0~;~2].

On note $h'$ sa dérivée et $h''$ sa dérivée seconde.

On admet que, pour tout réel $x$ de l'intervalle ]0~;~2], on a : 

\[h'(x) = 4x(1 + \ln (x)).\]

\begin{enumerate}[resume]
\item Sur l'intervalle $\left]\dfrac{1}{\text{e}}~;~2\right]$, la fonction $h$ s'annule :

\begin{center}
\begin{tabularx}{\linewidth}{*{2}{X}}
\textbf{a.~~} exactement 0 fois. &\textbf{b.~~}  exactement 1 fois.\\
\textbf{c.~~} exactement 2 fois. &\textbf{d.~~}  exactement 3 fois.
\end{tabularx}
\end{center}

\item Une équation de la tangente à $\mathcal{C}_h$ au point d'abscisse $\sqrt{\text{e}}$ est:

\begin{center}
\begin{tabularx}{\linewidth}{*{2}{X}}
\textbf{a.~~} $y = \left(6\text{e}^{\frac12}\right) .{} x $&
\textbf{b.~~} $y = \left(6\sqrt{\text{e}}\right).{} x + 2\text{e}$\\
\textbf{c.~~} $y = 6\text{e}^{\frac{x}{2}}$&\textbf{d.~~} $y = \left(6\text{e}^{\frac12}\right) .{} x - 4\text{e}$.
\end{tabularx}
\end{center}

\item Sur l'intervalle ]0~;~2], le nombre de points d'inflexion de la courbe $\mathcal{C}_h$ est égal à :

\begin{center}
\begin{tabularx}{\linewidth}{*{4}{X}}
\textbf{a.~~} 0&\textbf{b.~~} 1 &\textbf{c.~~} 2 &\textbf{d.~~}  3
\end{tabularx}
\end{center}

\item \footnote{Uniquement au Liban} On considère la suite $\left(u_n\right)$ définie pour tout entier naturel $n$ par

\[u_{n+1} = \dfrac12u_n + 3\quad \text{et}\quad u_0 = 6.\]

On peut affirmer que :

\begin{center}
\begin{tabularx}{\linewidth}{*{2}{X}}
\textbf{a.~~} la suite $\left(u_n\right)$ est strictement croissante.&\textbf{b.~~} la suite $\left(u_n\right)$ est strictement décroissante.\\
\textbf{c.~~} la suite $\left(u_n\right)$ n'est pas monotone. &\textbf{d.~~} la suite $\left(u_n\right)$ est constante.
\end{tabularx}
\end{center}
\end{enumerate}

\bigskip

\end{document}
